\section{REVIEW PROCESS AND EVIDENCE BASE}
\label{sec:review_methods}
\label{sec:corpus_appraisal_results}

With these comparison conditions established, the analysis turns to the
evidence base for the technical synthesis. We assembled the corpus through a
systematic review reported in accordance with the Preferred Reporting Items for
Systematic Reviews and Meta-Analyses (PRISMA) 2020 statement. Search reporting
followed the PRISMA extension for searches (PRISMA-S)
\cite{PRISMA2020,PRISMAS2021}. The following subsections describe corpus
construction, evidence extraction, and technical appraisal.
Table~\ref{tab:evidence_reconciliation} maps the analytical role of each
evidence category.

\begin{table*}[t]
\caption{Use of extracted evidence in the survey. Row counts represent coding
records, and study counts may overlap across categories.}
\label{tab:evidence_reconciliation}
\centering
\footnotesize
\setlength{\tabcolsep}{5.0pt}
\renewcommand{\arraystretch}{1.04}
\begin{tabularx}{\textwidth}{@{}>{\raggedright\arraybackslash}p{0.22\textwidth}
>{\raggedleft\arraybackslash}p{0.08\textwidth}
>{\raggedleft\arraybackslash}p{0.09\textwidth}
>{\raggedright\arraybackslash}X@{}}
\toprule
Evidence use & Rows & Studies & Role in the survey \\
\midrule
Quantitative evidence & 4,997 & 206 & Numerical or directional synthesis under recorded conditions. \\
Qualitative evidence & 3,206 & 206 & Mechanisms, trends, and limitations. \\
Claims used for context & 31 & 21 & Contextual interpretation reported
separately from the primary synthesis. \\
Claims with source conflicts & 72 & 31 & Recorded separately under predefined
decision rules, while supported evidence from the same studies remained in the
primary synthesis. \\
Primary synthesis & 8,203 & 206 & Quantitative and qualitative evidence combined. \\
\bottomrule
\end{tabularx}
\vspace{1pt}
\parbox{\textwidth}{\footnotesize\emph{Note.} The complete evidence register
contains
8,306 rows. Electronic Supplement S-Evidence reports the 8,203 primary
synthesis rows and their source locations. Electronic Supplement S-Protocol
documents the decision rules.}
\end{table*}

\subsection{Corpus Construction}

To assemble that evidence base, we searched Scopus, IEEE Xplore, ScienceDirect,
SpringerLink, Wiley Online Library, and Taylor \& Francis Online. The search
covered journal articles and full conference papers that had undergone peer
review and were first made publicly available from 1 January 2020 through
22 June 2026. Eligible reports contained sufficient technical content in
English, examined communication and sensing within an optical or photonic
system, and provided assessable evidence for at least one survey domain.
Reviews and other contextual sources supported field positioning. Eligible
primary reports proceeded to study reconciliation for the technical
denominator. Search strategies, eligibility decisions, and protocol
documentation are provided in Electronic Supplements S-Search, S-Exclusions,
and S-Protocol.

After deduplication, screening, and assessment of full texts, 227 reports met the
eligibility criteria. Among them, 21 were companion reports for studies also
represented by another eligible report. Consolidating eligible reports by study
produced 206 unique studies for the technical synthesis. The totals differ by
21 because companion reports map to studies already represented in the eligible
set. This mapping preserves provenance for every report while counting each
underlying study once. Electronic Supplement ST-01 cites all 206 studies and
traces their 227 reports.

Figure~\ref{fig:prisma_report_study_flow} reconciles the selection flow while
keeping records, reports, and studies as distinct counting units.

\begin{figure*}[!t]
\centering
\includegraphics[width=\textwidth]{figures/final_submission_2026-08-25/fig03_prisma_report_study_flow.pdf}
\caption{Selection of the O-ISAC evidence base. Six sources yielded 1,733
records and 1,259 entered title and abstract screening. The retrieval route
contained 330 unique reports after two bibliographic aliases were consolidated.
Assessment of full texts covered 272 reports and produced 227 eligible reports.
Companion mapping linked these reports to 206 unique studies. The contextual
corpus contains 67 records that support field positioning.}
\label{fig:prisma_report_study_flow}
\end{figure*}

The counting unit changes from records to reports at retrieval and from reports
to studies at synthesis. Contextual records support field positioning, while
the study set defines the technical denominator.

\subsection{Evidence Extraction and Use}

With the study denominator established, we extracted evidence at three linked
levels covering the study, report, and individual claim. Each result retained
its definition, unit, measurement plane, operating conditions, validation
setting, and exact source location. Records remained separate when conditions
differed. Each field contained a reported value or a defined availability
state.
Quantitative extraction used values stated by the source, while graphical
findings remained in their reported form. Electronic Supplement S-Data
Dictionary provides the complete coding schema.

Claims entered the analysis according to their role. Quantitative and
qualitative claims supported the primary synthesis within their recorded
contexts. Contextual claims and claims with source conflicts were recorded
separately. This structure preserved supported evidence from all 206 studies,
while comparisons across studies followed the alignment rules in
Table~\ref{tab:comparison_record}.

\subsection{Technical Appraisal and Synthesis}

To clarify how each study supports the synthesis, we assessed the technical
evidence and reporting of all 206 studies with the Technical Quality Assessment
Framework (TQAF) developed for this review. The framework applies deterministic
rules across eight dimensions. A separate overall contribution category
summarizes the support each study provides to this synthesis. The profile
indicates how evidence can support interpretation in its source setting,
comparison through a shared benchmark, and comparison across platforms.
Electronic Supplements S-Protocol and S-Appraisal provide the framework's
methodological scope, scoring rules, and records for each study.

Figure~\ref{fig:tqaf_profile} places the eight appraisal dimensions beside the
separate overall evidence contribution category. The profile distinguishes
support for source interpretation from support for transfer and reuse.

\begin{figure*}[!t]
\centering
\includegraphics[width=\textwidth]{figures/final_submission_2026-08-25/fig04_tqaf_profile.pdf}
\caption{Technical appraisal profile of the 206 included studies. Each row
assigns all studies to strong, adequate, or low evidentiary support under the
TQAF rules developed for this review. The interpretation group covers technical
relevance, metric clarity, reporting completeness, and limitation transparency.
The transfer and reuse group covers validation maturity, reproducibility,
benchmark readiness, and comparison admissibility. Overall evidence
contribution remains a separate summary category.}
\label{fig:tqaf_profile}
\end{figure*}

Reporting completeness and metric clarity support interpretation within the
source setting. Validation, reproducibility, benchmark readiness, and
comparison admissibility introduce further conditions for transfer across
studies. These distinct roles explain the profile and guide the structured
synthesis.

The structured narrative synthesis retained each result's native task and unit
together with its operating conditions and validation setting. This structure
preserved the meaning of heterogeneous results while allowing recurring
mechanisms and design relationships to be compared across the corpus. The
evidence was organized into 115 bodies across seven survey domains, connecting
synthesis conclusions to their supporting studies. Electronic Supplement
S-Bodies provides their criteria, ratings, and memberships. With this evidence
structure in place, the technical synthesis begins with the physical role of
optics and the mechanism that couples communication and sensing.
